\chapter{Pronunciation}

\begin{sourcebox}
\centering\itshape
There is little profit in silent reading without the ability to read aloud
\end{sourcebox}

\section{Importance of pronunciation}

We assume throughout that the reader is studying alone and that he merely wishes to read the language, not to speak it. Nevertheless, there are many reasons why he must learn to pronounce Russian words aloud, instantaneously and with ease. Otherwise he will be unable to read silently without being overly disturbed by inward hesitations; and in any case the ability to pronounce words and phrases aloud is an indispensable aid to remembering them. At the beginning of his course the student must practice pronunciation, no matter how eager he may be to proceed to the mathematical passages later in the book. He should keep in mind the practical principle adopted in this book that he will get relatively little profit from silent reading (or rather, decipherment) as long as he cannot read aloud.\corrnote{底本此处写作 ``He must accept the dictum of all modern linguists ...''。把这一教学主张说成“所有现代语言学家的定论”过于绝对；正文仅将主语范围收窄为本书采用的教学原则，其余措辞保留。}

Moreover, he should aim at a fairly correct pronunciation, since a bad one, if unhesitating and consistent, is no easier to learn. In the present book we have distinguished certain features of Russian pronunciation that are to be regarded as indispensable, and certain other features (some of them relegated here to smaller type) which, though equally part of standard pronunciation, can probably be neglected without much practical harm. The student who is working alone is advised to read through all the exercises aloud at a fairly rapid pace. Although silent reading is his final goal, he will merely deceive himself if he reads silently from the beginning.

\section{The six Russian vowel-sounds}

Five of the six Russian vowel-sounds distinguished for the purposes of this book will give no trouble, but the sound of \ru{ы} (No.~29) does not occur in English. The tongue is drawn back almost as far as for the 
\emph{oo}-sound in \emph{ooze}, and the lips are straightened out as for the \emph{ee}-sound in \emph{ease}, a difficult combination for most of us, since in English there is no back vowel without lip-rounding.\corrnote{底本把这里的“六个元音”直接作为唯一音系分析。现代分析对 \ru{и/ы} 的音位地位并不完全一致；为忠实保留原有教学体系，正文只加 ``distinguished for the purposes of this book'' 这一限定。}

If the reader is studying alone, it will perhaps be best for him to accept the familiar English \emph{i}-sound in a word like \emph{bill}, although the tip of the tongue should be a little further back in the mouth. He must experiment here, as with all the sounds in this chapter, aiming at consistency and facility; he himself must decide on the pronunciation he prefers and must then practice it. Help from a speaker of Russian is highly advantageous but not indispensable.

\section{Monosyllables for practice in pronunciation}

To illustrate the six vowel sounds we give the following list of monosyllables (Exercise 2.1). These words should not be scorned on the ground that most of them are nonmathematical, since we must begin with monosyllables containing a wide range of vowels and consonants (for the special features of Russian polysyllables see §12).\corrnote{底本此处误指 ``see §11''；本章多音节词与重音实际为 §12，故直接改正交叉引用。} The English meanings, unimportant in themselves from the mathematical point of view, are to be utilized in the following way. When a Russian word is pronounced aloud, the student should note its meaning from the English and then energetically visualize the meaning quite apart from the word by which it is expressed in English; when he pronounces \ru{бак} tank (see Ex. 2.1) he should at the same time visualize an actual tank, and the habit thus formed should be continued for the mathematical terms occurring throughout the book. Visualization of the meaning, a practice valuable in itself and especially congenial to mathematicians, is in fact the only means whereby the lone student can avoid excessive dependence on English words. For fluent reading in a foreign language it is obviously necessary to escape as far as possible from one's native tongue.

Some students may at first be unwilling to pronounce aloud \emph{all} the Russian words given in this chapter (though it is not a time-consuming task), on the ground that they ``do not intend to speak the language.'' Such students should reflect on the fact that they themselves often say ``reading Russian mathematics would be much easier if the Cyrillic letters did not `look so strange'.'' The only practical method proposed here for removing this admittedly serious (though temporary) obstacle is to read words aloud until the response of the tongue to the sight of the letters becomes as automatic as the response of a typist's fingers; for then the letters will no longer \emph{look} strange.\corrnote{底本写 ``The only practical method''，这是过强的教学断言；仅改为 ``the only practical method proposed here''，不改其余论述。} The words should be pronounced as loudly as circumstances will permit, and where silence is necessary, one should listen carefully to one's own ``inner voice''. The column of MR transliterations given with this exercise should be covered with a sheet of paper and consulted only after the Russian word has been pronounced aloud. Keep in mind that Russian pronunciation is much easier than English in the sense that, in general, each letter has a comparatively regular value.\corrnote{底本作 ``each letter is always pronounced in the same way''。俄语拼读规则确实较英语规则，但重音、元音弱化、清浊同化及软硬音等使“总是同一读音”不成立，故只改动这一最小范围。} (The relatively few exceptions are dealt with in the following sections.) For example, the Cyrillic letter \ru{у} (No.~21) normally has the sound of \emph{-oo-} in \emph{too} or \emph{toot}, never as in \emph{took}; thus the Mathematical Reviews transcription \emph{mul} (for \ru{мул} mule) indicates the pronunciation \emph{mool}; the transcription \emph{sup} indicates \emph{soop} (as in the English word \emph{soup}), \emph{tut} is \emph{toot}, etc.

\begin{keybox}[Exercise 2.1]
Pronounce the following monosyllables aloud.
\end{keybox}

\begingroup
\footnotesize
\setlength{\tabcolsep}{3.5pt}
\renewcommand{\arraystretch}{1.04}
\begin{longtable}{@{}>{\cyrillicfont}p{17mm} p{22mm} p{23mm} @{\hspace{4mm}} >{\cyrillicfont}p{17mm} p{22mm} p{23mm}@{}}
\toprule
\normalfont Russian & English & MR transcription & \normalfont Russian & English & MR transcription\\
\midrule
\endfirsthead
\toprule
\normalfont Russian & English & MR transcription & \normalfont Russian & English & MR transcription\\
\midrule
\endhead
а & and & a & дым & smoke & dym\\
акт & act & akt & жар & heat & žar\\
бак & tank & bak & жох & swindler & žoh\\
бал & ball & bal & жук & beetle & žuk\\
бант & bow & bant & за & beyond & za\\
бар & bar & bar & звон & peal & zvon\\
бас & bass & bas & звук & sound & zvuk\\
бок & side & bok & знак & sign & znak\\
бор & forest & bor & зонт & umbrella & zont\\
борт & side & bort & и & and & i\\
борщ & beet soup & boršč & ил & silt & il\\
бот & bot & bot & как & how & kak\\
брат & brother & brat & карп & carp & karp\\
брус & beam & brus & класс & class & klass\\
бы & would & by & кот & tomcat & kot\\
бык & bull & byk & Крым & Crimea & Krym\\
был & was & byl & курс & course & kurs\\
быт & life & byt & куст & bush & kust\\
вал & billow & val & ли & whether & li\\
ваш & your & vaš & лук & onion & luk\\
вол & ox & vol & луч & ray & luč\\
волк & wolf & volk & мак & poppy & mak\\
вон & there & von & Март & March & Mart\\
вот & here & vot & мол & pier & mol\\
врач & doctor & vrač & мост & bridge & most\\
гол & goal & gol & мул & mule & mul\\
гул & din & gul & мы & we & my\\
да & yes & da & мыл & washed & myl\\
дал & he gave & dal & мыс & cape & mys\\
дар & gift & dar & на & on & na\\
дом & house & dom & нас & us & nas\\
дух & spirit & duh & наш & our & naš\\
но & but & no & спорт & sport & sport\\
нос & nose & nos & стан & stature & stan\\
о & about & o & сто & hundred & sto\\
он & he & on & стол & table & stol\\
от & from & ot & сук & bough & suk\\
пакт & pact & pakt & суп & soup & sup\\
пар & steam & par & сын & son & syn\\
пат & stalemate & pat & сыр & cheese & syr\\
план & plan & plan & та & that & ta\\
плац & parade-ground & plac & так & thus & tak\\
плач & weeping & plač & там & there & tam\\
плащ & raincoat & plašč & танк & tank & tank\\
плуг & plough & plug & тип & type & tip\\
плут & rascal & plut & то & therefore & to\\
по & according to & po & ток & current & tok\\
пол & floor & pol & тост & toast & tost\\
порт & port & port & тот & this & tot\\
пост & fast & post & труп & corpse & trup\\
пот & perspiration & pot & трус & coward & trus\\
при & near & pri & тур & round & tur\\
прут & twig & prut & тут & here & tut\\
прыщ & pimple & pryšč & туш & flourish & tuš\\
пук & bundle & puk & ты & thou & ty\\
пункт & item & punkt & тыл & rear & tyl\\
пуск & starting & pusk & у & at & u\\
пуф & puff & puf & ум & mind & um\\
пух & eiderdown & puh & факт & fact & fakt\\
пыл & ardor & pyl & флот & fleet & flot\\
рак & cancer & rak & фон & background & fon\\
рост & growth & rost & фронт & front & front\\
рот & mouth & rot & хор & choir & hor\\
сам & self & sam & час & hour & čas\\
со & with & so & шар & sphere & šar\\
сок & juice & sok & шарф & scarf & šarf\\
сон & sleep & son & шкаф & cupboard & škaf\\
сор & rubbish & sor & штык & bayonet & štyk\\
 & & & шум & noise & šum\\
\bottomrule
\end{longtable}
\endgroup

\section{Preliminary remarks on hard and soft consonants}

For an English-speaking reader the chief difficulty in Russian pronunciation lies in the fact that the Russian hard--soft consonant contrast occurs in a wider variety of positions and functions more systematically than the nearest English approximations.\corrnote{底本写 ``although soft consonants are found in both languages''。英语中常用来类比的 consonant+$j$ 序列并不等同于俄语的音位性腭化；这里仅修正这一比较的精确性。}

As we have seen in Chapter I, three of the 21 Russian consonants are normally hard (\ru{ц, ж, ш}; Nos.~24, 8, 26) and three of them (\ru{й, ч, щ}; Nos.~11, 25, 27) are normally soft. For the other fifteen consonants, \ru{б, п, в, ф, д, м, н, р, л, т, з, с, г, к, х}, the hard pronunciation may be regarded as the normal one; and then we must ask: how is a hard consonant modified by being made soft? We shall consider this question (in English approximations and in Russian) for consonants occurring before the vowels in each of the five columns of the vowel-scheme, and then elsewhere, i.e. before another consonant or at the end of a word (six cases in all).

\section{Hard and soft consonants in English and Russian}

\subsection*{Case i: before the \emph{u}-sound (first column of the vowel-scheme)}

Here the contrast between hard and soft consonants is very noticeable in Russian and can be approximated by the English pairs in the following Exercise 2.2(i). Note that in all the words (English and Russian) the vowel-sound is intended to be the same, namely the back \emph{oo} as in \emph{ooze}. Since it is natural for a Russian consonant to be hard before a back vowel, many of our examples for soft consonants are necessarily words of foreign origin (\ru{бюст} bust, \ru{пюре} purée, etc.).

The reader should pronounce these English and Russian words several times, noting for each pair the difference in the position of the tongue. Historically and articulatorily, Russian softness may suggest a relationship with a following palatal gesture, but in present-day Russian a soft consonant is a single palatalized consonant, not a hard consonant followed by a separate /j/. The English contrast between \emph{coo}, \emph{cue}, and \emph{seek-you} can therefore serve only as an approximation: most modern analyses treat the beginning of English \emph{cue} as /k/ followed by /j/, not as the Russian phoneme $/k^{j}/$.\corrnote{底本把 English \emph{cue} 中的 k 直接称为 ``a soft k''，并与俄语软辅音等同；现代英语音系通常分析为 /k+j/。保留原有教学类比，但直接改正其音系身份。} Thus the usual transcription of the Russian word \ru{тюк} package as \emph{tyuk} may suggest an inaccurate pronunciation (t followed by consonantal y), whereas the intended Russian sound is a single palatalized consonant. Nevertheless, we shall continue to use this transcription for its typographical convenience. At the end of a word (or before another consonant) we shall use a prime, as in Mathematical Reviews, e.g. \ru{брать} (\emph{brat'}) take.

\begin{keybox}[Exercise 2.2(i) --- hard and soft consonants before /u/]
\textit{English mnemonic pairs:} booty--beauty; pooh--pew; voodoo--view; fool--few; moo--mew; noon--new; room--(Ryukyu?); loom--lewd; do--dew; too--tube; zoo--Zeus; soup--sue; goo--ague; kook--Kew; hoot--hew.

\medskip
\textit{Russian pairs.} The accent is marked on the few words of more than one syllable.
\end{keybox}

\begingroup
\small
\setlength{\tabcolsep}{4pt}
\renewcommand{\arraystretch}{1.03}
\begin{longtable}{@{}>{\cyrillicfont}p{29mm} p{38mm} >{\cyrillicfont}p{34mm} p{40mm}@{}}
\toprule
\normalfont Hard & Meaning & \normalfont Soft & Meaning\\
\midrule
\endfirsthead
\toprule
\normalfont Hard & Meaning & \normalfont Soft & Meaning\\
\midrule
\endhead
бук & beech & бюст & bust\\
пук & bundle & пюре́ & purée\\
вуз & institute for higher education & ревю́ & revue\\
фуру́нкул & furuncle & фюзеля́ж & fuselage\\
му́зыка & music & мюзик-холл & music-hall\\
ну & now & нюх & flair\\
рука́ & hand & рюмка & wine-glass\\
луч & ray & люк & hatchway\\
душ & shower-bath & дю́на & dune\\
стук & rap & тюк & package\\
Зу́лу & Zulu & изю́м & raisins\\
сук & bough & сюда́ & hither\\
гул & rumbling & Гю́го & (Victor) Hugo\\
кула́к & kulak & кюри́ & curie\\
худо & evil & Хуэ́ & Huế (name of the city)\\
\bottomrule
\end{longtable}
\endgroup

Note that \emph{room} has no precise companion in English; consequently, many speakers of English will make no distinction between hard and soft \ru{р} in Russian; e.g. they will pronounce the first syllable of \ru{рюмка} wine-glass in whatever way they pronounce the English word \emph{room}. (Actually the Russian \ru{р}, No.~18, is normally trilled or tapped; and the reader should make the trill, if it is natural for him.)

\subsection*{Case ii: before the \emph{o}-sound (second column of the vowel-scheme)}

In English it is hard to find exact examples of Russian-style soft consonants before the \emph{o}-sound (which in Russian will here be written \ru{ё} as in \ru{Фёдор} Fyodor, the given name of the novelist Dostoevski). But approximate examples can be manufactured by pronouncing words like \emph{junior} and \emph{senior} as though they were \emph{june-yor} and \emph{seen-yor}; these remain mnemonic approximations rather than exact phonetic equivalents.\corrnote{底本直接把这些英语序列当作软辅音实例；这里仅注明其近似性质。}

As is to be expected, most consonants are hard before this back \emph{o}-sound, a fact which means, as we have seen, that the letter \ru{ё} (No.~7) is relatively infrequent. When it does occur, it normally bears the word stress, and in printed Russian texts the ``double dot'' is often omitted (though we shall retain it in this book). The omission does not make \ru{е} and \ru{ё} the same letter, and modern dictionaries differ in how systematically they print \ru{ё}.\corrnote{底本称省略两点的结果是 \ru{е} 与 \ru{ё} 在俄语词典排序中“不加区分”。现代辞书实践并不完全统一，且 \ru{ё} 仍是独立字母，故作最小修正。} Here are four pairs of Russian words illustrating the difference between \ru{о} (No.~16) and \ru{ё} (No.~7).

\begin{keybox}[Exercise 2.2(ii)]
\begin{tabular}{@{}>{\cyrillicfont}l l @{\hspace{4mm}} >{\cyrillicfont}l l@{}}
дом & house & идёт (i-dyot) & goes\\
пот & perspiration & пёс (pyos) & dog\\
торг & trade & тётя (tyotya) & aunt\\
фок & foresail & Фёдор (Fyodor) & Theodore\\
\end{tabular}
\end{keybox}

\subsection*{Case iii: before the \emph{a}-sound (third column of the vowel-scheme)}

Here again it is hard to find exact Russian-style soft consonants in English, but we can get approximate practice by slightly mispronouncing the names of flowers \emph{dahlia}, etc., or of girls \emph{Claudia}, etc., or of countries \emph{Ethiopia}, etc., as though they were spelled as follows: \emph{dah-lyah}, \emph{bego-nyah}, \emph{sal-vyah}; \emph{Clau-dyah}, \emph{Zeno-byah}; \emph{Ethio-pyah}, \emph{Bohe-myah}, etc.\corrnote{这些英语例子只提供发音近似，并非与俄语腭化辅音完全相同。}

Here are seven pairs of Russian words illustrating the difference between \ru{а} (No.~1) and \ru{я} (No.~34).

\begin{keybox}[Exercise 2.2(iii)]
\begin{tabular}{@{}>{\cyrillicfont}l l @{\hspace{4mm}} >{\cyrillicfont}l l@{}}
вал & billow & вял (vyal) & flabby\\
да & yes & дядя & uncle\\
зам- & vice- (e.g. president) & взят & taken\\
мак & poppy & мяч & ball\\
на & on & няня & nurse\\
сак & sack & сяк & thus\\
так & thus & тятя & daddy\\
\end{tabular}
\end{keybox}

\subsection*{Case iv: before the \emph{e}-sound (fourth column of the vowel-scheme)}

An approximate English example is provided by the word \emph{lorgnette} (\emph{lornyet}), the last sounds of which may serve as a mnemonic for the Russian word \ru{нет} no, familiar to American newspaper readers in the approximative spelling \emph{nyet}.

\begin{keybox}[Exercise 2.2(iv)]
\begin{tabular}{@{}>{\cyrillicfont}l l @{\hspace{4mm}} >{\cyrillicfont}l l@{}}
век (vyek) & century & нет & no\\
лес & forest & пел & sang\\
мел & chalk & сел & sat\\
\end{tabular}
\end{keybox}

The letter \ru{э} is also common in transcriptions of non-Russian names and in loanwords, e.g. \ru{Бэла} in Lermontov's \emph{A Hero of Our Time} and \ru{Брауэр} for Brouwer.

In native Russian words, the front \emph{e}-sound normally softens an immediately preceding paired consonant, which means that after such a consonant it is ordinarily written \ru{е} rather than \ru{э}. The letter \ru{э} is rare in inherited Russian vocabulary but is not confined to a single word; it occurs in forms of \ru{этот} and in a small number of other native or expressive items.\corrnote{底本写 ``the letter \ru{э} does not occur at all in any native Russian word (except \ru{этот})''。这过于绝对，故直接改正为“固有词中罕见，但不只一个词”。}

But \ru{э} may occur freely in foreign words. For example, \ru{Бэла} (pronounced like our Bella) is the name of the Georgian (i.e. non-Russian) heroine in Lermontov's novel \emph{A Hero of Our Time}; and the reader will encounter \ru{э} in transcriptions of names of non-Russian mathematicians, e.g. \ru{Брауэр} Brouwer.

\subsection*{Case v: before the sounds of \ru{и} (No.~10) and \ru{ы} (No.~29) (fifth column of the vowel-scheme)}

Here the soft consonants (before \ru{и}) will give the reader relatively little trouble since English approximations often encourage a forward tongue position; e.g. the word \emph{teak} can serve as a rough mnemonic for Russian \ru{тик}. Before \ru{ы}, paired consonants have the normal hard pronunciation.\corrnote{底本说 English \emph{teak} 与 Russian \ru{тик} “pronounced in the same way”。英语 /tiːk/ 中的 /t/ 与俄语 $/t^{j}/$ 并不相同，故只改为近似助记。}

\begin{keybox}[Exercise 2.2(v)]
\textit{Soft consonants before \ru{и}}
\end{keybox}

\begingroup
\small
\setlength{\tabcolsep}{1.5pt}
\begin{longtable}{@{}>{\cyrillicfont}p{20mm} p{23mm} >{\cyrillicfont}p{20mm} p{23mm} >{\cyrillicfont}p{20mm} p{26mm}@{}}
бил & struck & лик & image & писк & squeak\\
бинт & bandage & лист & leaf & Рим & Rome\\
вил & wove & лифт & elevator & риск & risk\\
винт & screw & мим & mime & сил & of forces\\
гипс & gypsum & мир & world & тик & teak\\
диск & disk & миф & myth & тип & type\\
кист & bunch & ни & neither & тис & yew\\
кит & whale & пик & peak & тиф & typhus\\
ли & whether & пир & feast & & \\
\end{longtable}
\endgroup

\begin{keybox}
\textit{Hard consonants before \ru{ы}}
\end{keybox}

\begingroup
\small
\setlength{\tabcolsep}{1.5pt}
\begin{longtable}{@{}>{\cyrillicfont}p{20mm} p{23mm} >{\cyrillicfont}p{20mm} p{23mm} >{\cyrillicfont}p{20mm} p{26mm}@{}}
бы & would & мыс & cape & сыр & cheese\\
был & was & нытик & whimperer & сыск & investigation\\
вы & you & пыж & wad & сыч & owl\\
дым & smoke & пыл & ardor & ты & thou\\
зыбка & cradle & пытка & torture & тыл & rear\\
лыжи & skis & рыба & fish & тын & paling\\
лыко & bast & рысак & trotter & & \\
мы & we & сын & son & & \\
\end{longtable}
\endgroup

\subsection*{Case vi: at the end of a word, or before another consonant}

For this case, in which softness of a consonant is often indicated by the soft sign \ru{ь} (No.~30), there is very little exact help to be found in English. But we can proceed approximately as follows. If the reader were asked to pronounce the initial consonantal sound of \emph{booty} he would unhesitatingly reply with the sound of a hard b. The beginning of \emph{beauty}, \emph{pew}, \emph{tube}, etc., may be used as a rough articulatory mnemonic for a Russian soft consonant, but in English these are normally consonant+/j/ sequences rather than isolated Russian-type palatalized phonemes.\corrnote{底本称 English \emph{beauty} 的起始辅音就是“soft b”，并建议直接抽离；这里保留练习意图但改正音系描述。} The student should practice the Russian soft consonants until he can easily pronounce them at the end of a word. The soft \ru{л} and \ru{н} may be approximated from English sequences such as those in \emph{million} and \emph{canyon}, though again the correspondence is not exact.

\begin{keybox}[Exercise 2.2(vi)]
\begin{tabular}{@{}>{\cyrillicfont}l l @{\hspace{4mm}} >{\cyrillicfont}l l @{\hspace{4mm}} >{\cyrillicfont}l l@{}}
бить & beat & киль & keel & сильф & sylph\\
быль & fact & линь & rope-line & соль & salt\\
вить & weave & лить & pour & фильм & film\\
дать & give & роль & role & нить & thread\\
день & day & & & &\\
\end{tabular}
\end{keybox}

\section{A first approximation to Russian pronunciation}

The suggestions given at the beginning of Chapter I for the pronunciation of the various letters of the alphabet are based on an extremely simple and concise, though only approximate system of describing Russian pronunciation; namely, each of the consonants is said to have only one pronunciation (in general, the one up to now called ``hard''); the hard and soft signs are said to be silent; and the vowels \ru{ю, ё, я, е}, which along with \ru{и} are now often called ``soft vowels'', are said to be pronounced like the corresponding \ru{у, о, а, э} with prefixed consonantal y (namely, yu, yo, ya, ye).

But the system is only a rough approximation. In particular, soft consonants at the end of a word remain palatalized rather than being pronounced as though hard, and all soft consonants are misrepresented if they are treated simply as hard consonants followed by /j/.\corrnote{底本原文反而说 ``soft consonants at the end of a word are pronounced as though they were hard''，这是明确错误；标准俄语的词末软辅音保持腭化，故直接改正。} In attempting to make the pronunciation easier for readers who have no interest in speaking Russian, the system gives up the distinction between hard and soft consonants, thereby obliterating one of the most noticeable and characteristic features of the language. The phrase ``soft vowel'' is therefore best understood only as a traditional pedagogical label for vowel letters which can indicate the softness of a preceding consonant; vowels themselves do not bear the Russian consonantal hard--soft contrast.\corrnote{底本用“辅音会产生 frictional noise”来定义 consonant，从而论证 ``soft vowel'' 是自相矛盾；这一语音学定义并不成立（如塞音并非持续摩擦音），故仅保留其结论性教学提醒并改正理由。} Nevertheless, in a textbook written for speakers of English, the system is convenient for preliminary use when the letters of the alphabet are first introduced.

\section{The letter \ru{й} when pronounced but not written}

In Chapter I we have seen that the four vowel-symbols \ru{ю, ё, я, е} may indicate a preceding soft consonant after a consonant letter, but when no consonant is written before them (and also after a vowel or separating sign) they normally include the sound /j/ before the vowel.\corrnote{底本写作这四个符号 ``can occur only after a soft-consonant sound''，并因此把词首 /j/ 解释为“自动理解出一个未写的软辅音 \ru{й}”。正字法上这不是唯一环境；词首、元音后和分隔符号后这些字母本身就表示 /j/+元音，故直接改正。} For the elementary reading convention used here, it is convenient to think of an unwritten \ru{й}-sound in such cases, as in the following exercise.

\begin{keybox}[Exercise 2.3]
\ru{ел} ``ate'' may be read as though it began with /j/: approximately \emph{yel}; similarly \ru{ёлка} ``spruce'', \ru{семья} ``family'', \ru{съел} ``ate up'', \ru{трение} ``friction'', \ru{юла} ``whirligig'', \ru{юмор} ``humor'', \ru{ют} ``quarter-deck'', \ru{я} ``I'', \ru{яма} ``hole''.
\end{keybox}

\section{The ``separating'' hard and soft signs}

Since a consonant is hard unless the spelling specifically indicates that it is soft, the hard sign \ru{ъ} (No.~28) might seem to be superfluous; but this is not quite the case. For example, in the preceding list the hard sign is necessary in \ru{съел} ate up, with the sounds s (hard), y, e, l, to distinguish it from \ru{сел} sat, in which the initial consonant is soft. If a prefix ending in a hard consonant, e.g. \ru{под-} under, \ru{из-} out of, \ru{об-} around, \ru{от-} from, \ru{с-} off, is prefixed to a word beginning with \ru{ю, ё, я} or \ru{е}, a so-called separating hard sign is inserted between them, e.g. \ru{объём} volume, in order to preserve the hardness of the final consonant in the prefix.

After many native prefixes ending in a consonant, word-initial \ru{и} changes orthographically to \ru{ы}; thus \ru{интеграл} integral but \ru{подынтегральный} under the integral sign. In some formations and especially where the morphological structure is to be kept transparent, \ru{и} may be retained.\corrnote{底本把这一 \ru{и$\to$ы} 现象与“separating hard sign”并列叙述；它实际上是另一条形态音系/正字法规则，故仅把关系说清。}

Similarly the separating soft sign in \ru{семья} family shows that there is an extra /j/-sound here in comparison with \ru{семя} seed while the \ru{м} remains soft.

These separating signs occur fairly often in mathematical texts and in names, e.g. \ru{Стилтьес} Stieltjes, \ru{Ньютон} Newton, etc. The present book may sometimes simplify them in approximate reading, but they are not phonologically meaningless and should not in general be ignored.\corrnote{底本结尾写 ``in practice it is convenient to ignore them''。分隔符号可决定 /j/ 是否出现以及前一辅音的软硬，故不能作为一般规则忽略；这里只修正该过强建议。}

\section{The letter \ru{й} when written}

In native Russian words the letter \ru{й} is written especially after vowels and frequently at the end of a syllable; thus \ru{май} May, \ru{бей} strike, \ru{бой} battle, \ru{белый} white, \ru{высокий} high. It is not, however, restricted to the five preceding vowel letters listed in the original formulation, and modern Russian also uses initial \ru{й} in some borrowings and proper names.\corrnote{底本写 \ru{й} “only at the end of a syllable and only after the vowels \ru{а, е, о, ы, и}”，该限制并不成立，故作最小范围修正。} Modern phonetic descriptions generally treat the sound as the consonantal approximant /j/; ``semivowel'' is an alternative descriptive term rather than an unresolved physical question.\corrnote{底本说俄语语言学家对于 \ru{й} 是 consonant 还是 semivowel 的争议“is a question to be settled in a physics laboratory”。现代语音学通常把它列为近音辅音 /j/，故直接更新这一事实判断。}

Here we shall merely advise the reader to continue calling \ru{й} a consonant. The sequences \ru{ай, ей, ой} may be pronounced with a final [j]-glide and can be compared approximately with the English sequences in \emph{aisle}, \emph{eight}, \emph{oil}; but the \ru{й} in adjective endings such as \ru{-ый} and \ru{-ий} is not altogether silent in standard Russian.\corrnote{底本建议把 \ru{-ый/-ий} 中的 \ru{й} ``altogether silent''；现代标准读音中通常仍有 [j] 或相近滑音成分，故直接改正。}

\section{Assimilation of voiced and voiceless consonants}

In the following situation a Russian consonant changes its pronunciation with no corresponding change in the spelling.

The twenty-one Russian consonant letters fall, for the purposes of this discussion, into voiced and voiceless classes. For a voiced consonant the vocal cords vibrate. The reader may experiment, say, with \ru{б} No.~2 (voiced) and \ru{п} No.~17 (voiceless), or \ru{д} No.~5 (voiced) and \ru{т} No.~20 (voiceless). Some consonants occur in pairs, voiced and voiceless (but otherwise similar), and others stand alone, according to the following table. The source table accidentally omitted \ru{р} from the voiced sonorants; it is restored here.\corrnote{底本标题称有21个辅音字母，但表中实际只列20个，漏列有声响音 \ru{р}；现直接补入。}

\begin{center}
\small
\begin{tabular}{@{}l*{11}{>{\cyrillicfont}c}@{}}
\toprule
Voiceless & ц & ч & щ & х & к & т & с & ш & п & ф & \\
Voiced    & г & д & з & ж & б & в & л & м & н & р & й\\
\bottomrule
\end{tabular}
\end{center}

When a voiced and voiceless obstruent occur together, in either order (e.g. \ru{отбор} selection or \ru{вправо} to the right), assimilation commonly makes their voicing agree. In English the situation is quite complicated. Assimilation may or may not occur; when it does occur, it may or may not affect the spelling; and its direction depends on the construction. For example, in the word \emph{roped} the written \emph{d} is pronounced voiceless (contrast \emph{robed}) by assimilation to the voiceless \emph{p}.

In Russian this voiced--voiceless assimilation behaves much more regularly. In the principal cases considered here it is regressive (the later obstruent conditions the preceding one), although detailed phonetic rules also depend on the type of consonant and word boundary.\corrnote{底本把“第一个辅音同化于第二个辅音”称为 ``forwards''；现代术语中这正是逆行/回归同化（regressive assimilation），故直接改正术语。} The practical rules are:

\begin{enumerate}[label=\alph*)]
\item before a voiceless obstruent, or at the end of a phonological word, a paired voiced obstruent is normally devoiced, i.e. pronounced like its voiceless partner;
\item before the paired voiced obstruents \ru{г, д, з, ж, б}, the corresponding paired voiceless obstruents \ru{к, т, с, ш, п} normally become voiced. The consonant \ru{в} has special behavior and ordinarily does not trigger the same voicing of a preceding obstruent.
\end{enumerate}

In connected speech most of these changes will be made automatically, and in fact, if the student ignores some of them, relatively little harm will be done, provided he is still able to read without hesitation.

\begin{keybox}[Exercise 2.4]
Pronounce aloud:
\end{keybox}

\begingroup
\small
\setlength{\tabcolsep}{1.5pt}
\begin{longtable}{@{}>{\cyrillicfont}p{22mm} p{17mm} p{28mm} >{\cyrillicfont}p{22mm} p{17mm} p{28mm}@{}}
\multicolumn{6}{@{}l}{\textit{a) Devoicing}}\\[2pt]
все & (fsye) & everybody & миг & (mik) & moment\\
вход & (fhot) & entrance & раз & (ras) & time\\
входить & (fhodit') & enter & резка & (reska) & cut\\
дуб & (dup) & oak & рог & (rok) & horn\\
дубка & (dupka) & small oak & сад & (sat) & garden\\
лоб & (lop) & forehead & хлеб & (hlep) & bread\\
лов & (lof) & catch & юг & (yuk) & south\\
лодка & (lotka) & boat & яд & (yat) & poison\\[4pt]
\multicolumn{6}{@{}l}{\textit{b) Voicing and cluster assimilation}}\\[2pt]
книжка & (knižka) & booklet & сгиб & (zgip) & bend\\
ловко & (lofko) & cleverly & сказка & (skaska) & tale\\
робкий & (ropki) & timid & сделать & (zdelat') & to do\\
сверх & (sverh) & beyond & экзамен & (egzamen) & examination\\
\end{longtable}
\endgroup

\section{Consonant clusters}

The clusters of consonants permitted in given settings, especially at the beginning or end of a word, are highly characteristic of a given language. For example, most foreigners have difficulty with the English word \emph{strengths}, and Russian has many combinations not found in English.

In the very common word \ru{много} much, many (e.g. in \ru{многочлен} = many-term = polynomial, imitated from the Greek), the reader should first pronounce a protracted \ru{м} (a humming sound), and then add \ru{-ного} to it, trying to avoid any impression that the \ru{м} forms an extra syllable; in other words, he should proceed in exactly the same way as with the s in an English word like \emph{squelched}. This principle may be applied to all initial consonants that have a continuable sound; i.e. all those in Exercise 2.5 except \ru{г, д, к, п, т}. The most troublesome words are those with two non-continuable sounds, one after the other, e.g. as in \ru{где} where, for which the tongue, after first pronouncing the \ru{г}, should move forward for the \ru{-де} as smoothly as possible, with the idea of making the whole word into a monosyllable. Note also that certain Russian prepositions, e.g. \ru{в} in, consist of a single consonant, so that \ru{в доме} in the house is pronounced in connected speech with an initial consonant cluster.

Initial combinations of this sort are common in many languages; for example, in the Greek words from which we have derived \emph{pneumonia, pseudonym, psychology, pterodactyl, gnostic}, etc. When speaking English, we nowadays refuse to pronounce the first consonant, but this does not apply to the corresponding Russian words \ru{пневмония, псевдоним}, etc. However, when a consonant is doubled in Russian, its phonetic realization is not universally identical to a single consonant: in careful speech some double consonants are long, while in other positions or styles they may be shortened.\corrnote{底本写 ``when a given consonant is doubled in Russian, it is pronounced as though single''，这一概括过强；俄语双辅音的实际时长取决于词位、语素边界与语体，不能一概等同单辅音。} Usually, only one of the consonants is written in Russian when the morphology and orthography require a single letter; e.g. \ru{агрегат} aggregate.

\begin{keybox}[Exercise 2.5]
Pronounce aloud these words containing consonant clusters.
\end{keybox}

\begingroup
\small
\setlength{\tabcolsep}{1.5pt}
\begin{longtable}{@{}>{\cyrillicfont}p{21mm} p{16mm} p{29mm} >{\cyrillicfont}p{21mm} p{16mm} p{29mm}@{}}
вдаль & & far & вход & (fhot) & entrance\\
взад & (vzat) & backwards & гнать & & drive\\
взгляд & (vzglyat) & look & гнев & & anger\\
вздох & & deep breath & гнёт & (gnyot) & yoke\\
взлом & & breaking in & два & & two\\
взмах & & stroke & для & & for\\
взрыть & & plough & звать & & call\\
вклад & (fklat) & deposit & здравие & & health\\
вне & & outside & зло & & evil\\
всё & (fsyo) & all & змей & & serpent\\
знак & & sign & сдать & (zdat') & be weakened\\
квас & & kvass & средь & & among\\
кнут & & knout & срыв & (sryf) & frustration\\
кто & & who & счесть & & count\\
лгун & & liar & счёт & & calculation\\
мгла & & haze & тварь & & creature\\
мне & & to me & ткань & & cloth\\
много & & much & храм & & temple\\
мстить & & avenge & цвет & & color\\
нрав & (nraf) & disposition & член & & member\\
птица & & bird & чрез & & across\\
рвать & & to tear & шкаф & & cupboard\\
сбить & (zbit) & knock down & шнур & & cord\\
сбыт & (zbyt) & sale & штаб & (shtap) & staff\\
свет & & light & & &\\
\end{longtable}
\endgroup

Note that the common conjunction \ru{что} ``that'' is normally pronounced approximately \emph{shto} in standard Russian; spelling and pronunciation differ here.

Note that the very common conjunction \ru{что} = that has the irregular standard pronunciation approximately \emph{shto}.

\section{Words of more than one syllable; accents}

English is the monosyllabic language par excellence. In no other language is it possible to find pages of mathematics on which almost all the words are monosyllables, e.g. ``at this point the two curves have the same slope'', etc. In Russian the situation is far different. For example, in the standard translation of Jane Austen's novel \emph{Pride and Prejudice}, the translator finds it necessary to explain at some length why the hero is twitted by his friend for using so many four-syllable words. As we have seen, there are plenty of monosyllables in Russian, but the English-speaking reader must accustom himself to the fact that four-syllable words are just as common, and words of seven or eight syllables are not at all rare.

So the question arises: how are they to be pronounced?

In both English and Russian, polysyllabic words have an emphatic stress on the accented syllable (e.g. in English, \emph{possibility}) and the position of this accent (not given in ordinary Russian printed texts) is governed by complicated rules which we shall not try to discuss here (but recall that in Russian \ru{ё}, when written, normally carries the word stress). Since for the sake of pronouncing a word aloud it is important for the reader to know where the accent falls and since, on the other hand, we wish to have our texts look as natural as possible, we shall make the convention in this book that when the accent is left unmarked it is assumed to fall on the first syllable of the word (but all proper names are left unaccented). This is a convention of the present book, not a rule of Russian stress.\corrnote{底本的“未标重音即在第一音节”是教材内部约定；为避免被误读为语言事实，补明其性质。} As for the mathematical symbols and formulas, the practical solution is to read them as in English.

The basic difference emphasized here between English and Russian polysyllables is that a Russian simple word normally has one primary lexical stress; nevertheless, secondary stress can occur in compounds and some complex formations, so it is too strong to say that Russian has no secondary accents at all.\corrnote{底本写 ``in Russian there are no secondary accents'' 以及每个词 ``has exactly one accent''。这对于复合词及部分复杂派生词过于绝对，故作最小范围修正。} The reader should still make the primary stress an emphatic one, gliding lightly and rapidly over the unstressed syllables. All the preceding syllables may be regarded as preparation for the accent, and all the following ones as gliding away from it.

A long Russian word like \ru{непересекающиеся} non-intersecting (eight syllables) can be difficult for an English-speaking reader if all the unstressed syllables are given full, citation-form vowel qualities. Consequently, the gliding process is helped by the following useful but incomplete rule:

\begin{sourcebox}
In syllables preceding the accent, unstressed \ru{я} and \ru{е} after soft consonants are often reduced toward an \ru{и}-like quality.
\end{sourcebox}

Thus the practical pronunciation of \ru{непересекающиеся} can be approximated by the transcription used in the source, and similarly for forms such as \ru{спрямляемый} rectifiable. This is only an incomplete rule: standard Russian unstressed-vowel reduction depends on consonant hardness/softness, position relative to stress, and speech style.\corrnote{底本把“重音前 \ru{я、е} 都读成 \ru{и}”作为简化规则；现保留其教学用途，同时明确它是不完整的弱化规则。}

Various other rules for unstressed syllables, particularly with respect to the vowels \ru{а} and \ru{о}, are not developed here, although they form part of standard Russian. Northern dialects may show much less /o/ reduction than the Moscow-based standard, but it is not correct to say that the standard reductions are simply disregarded by ``many native speakers'' without qualification.\corrnote{底本把北方口音与标准语混为一谈；这里仅补上方言/标准语区分。}

The rules for unstressed vowels are not merely accidental results of rapid speech; they are systematic features of standard Russian phonetics, though their phonetic extent varies with speaking rate and style.\corrnote{底本写 ``all rules referring to the unstressed parts of words are merely the result of rapid speech''，这是明确的语音学错误，故直接改正。} The reader should nevertheless glide smoothly over unstressed syllables, keeping in mind that the Russian language, in spite of its numerous consonants, can in fact be spoken very rapidly and smoothly.

Here are some examples from the earlier part of the book:

\begin{center}
\begin{tabular}{@{}>{\cyrillicfont}l l l@{}}
определена́ & defined & o-pree-dee-lee-NA\\
произво́дная & derivative & pro-eez-VOD-na-ya\\
произведе́ние & product & pro-eez-vee-DYE-nee-ye\\
налега́ть & lie on & na-lee-GAT'\\
внесе́ние & insertion & vnee-SYE-nee-ye\\
\end{tabular}
\end{center}

These words are all of Class Three (see §2 of the Introduction), and therefore the exercises on their pronunciation will be postponed to Chapter V, where their formation and meaning will be discussed systematically; the following exercise contains only international words of Class Two.

\begin{keybox}[Exercise 2.6]
Pronounce each of these international words aloud, keeping in mind that the fundamental purpose of the exercise is to increase one's familiarity with the appearance of Cyrillic letters.
\end{keybox}

\begin{sloppypar}
\small
\setlength{\emergencystretch}{7em}
\textbf{a) Nouns ending in \ru{-а}:} e.g. \ru{аксио́ма} = axiom, \ru{амплиту́да} = amplitude: \ru{абсци́сса, арифме́тика, асимпто́та, брахистохро́на, гру́ппа, директри́са, координа́та, ло́гика, матема́тика, ма́трица, но́рма, резольве́нта, систе́ма, температу́ра, теоре́ма, фо́рмула, характери́стика, цикло́ида}.

\textbf{b) Nouns in \ru{-ция} ``-tion'':} \ru{абстра́кция, аксиоматиза́ция, аппроксима́ция, вариа́ция, девиа́ция, идентифика́ция, инду́кция, информа́ция, итера́ция, классифика́ция, компактифика́ция, конструкция, корреля́ция, линеариза́ция, минимиза́ция, опера́ция, оптимиза́ция, ортогонализа́ция, параметриза́ция, прое́кция, формализа́ция, фу́нкция}.

\textbf{c) Nouns in \ru{-ивность} ``-ivity'':} \ru{аддити́вность, ассоциати́вность, дистрибути́вность, коммутати́вность, рефлекси́вность, транзити́вность}. Cf. also \ru{-ность} ``-ity'': \ru{абсу́рдность, иррациона́льность, пропорциона́льность, сепара́бельность, стаби́льность}.

\textbf{d) Nouns in \ru{-ование}:} \ru{дифференци́рование, интегри́рование, сумми́рование}.

\textbf{e) Nouns ending in a hard consonant:} \ru{аппара́т, аргуме́нт, а́том, ба́зис, бикомпа́кт, ве́ктор, детермина́нт, диа́метр, дискримина́нт, дифференциа́л, идеа́л, изоморфи́зм, и́ндекс, интегра́л, интерва́л, класс, компоне́нт, конти́нуум, ко́смос, коэффицие́нт, ма́ксимум, меридиа́н, ме́тод, минимакс, ми́нимум, моме́нт, опера́тор, пара́метр, полиго́н, радика́л, сегме́нт, силлоги́зм, си́мплекс, спектр, факт, фа́ктор, факториа́л, фо́кус, фу́нктор, функционал, хара́ктер, центр, экстремум, элеме́нт, эллипсо́ид}.

\textbf{f) Adjectives in \ru{-ивный} ``-ive'':} \ru{аддити́вный, дедукти́вный, конструкти́вный, мультипликати́вный, примити́вный, рекурси́вный, рефлекси́вный, транзити́вный}.

\textbf{g) Adjectives in \ru{-ный}:} \ru{биголомо́рфный, гомеомо́рфный, изомо́рфный}.

\textbf{h) Adjectives in \ru{-альный}:} \ru{вертика́льный, дифференциа́льный, интегра́льный, инфинитезима́льный, лока́льный, минима́льный, натура́льный, норма́льный, пропорциона́льный, рациона́льный, спектра́льный, специа́льный, тривиа́льный, фундамента́льный}.

\textbf{i) Other adjectives in \ru{-ный}:} \ru{абсолю́тный, аналоги́чный, бина́рный, дискре́тный, коллинеа́рный, ко́мплексный, перпендикуля́рный, поля́рный, регуля́рный, сепара́бельный, симметри́чный, скаля́рный, трансценде́нтный, эквивале́нтный}.

\textbf{j) Adjectives in \ru{-ический} ``-ic/-ical'':} \ru{алгебраи́ческий, логи́ческий, геометри́ческий, гиперболи́ческий, канони́ческий, класси́ческий, логарифми́ческий, математи́ческий, метри́ческий, практи́ческий, тригонометри́ческий, цикли́ческий}.
\normalsize
\end{sloppypar}

\begin{keybox}[Exercise 2.7]
The names of these Russian mathematicians are to be pronounced aloud. The reader is invited to accent them in whatever way he may be accustomed to hearing them pronounced in English.
\end{keybox}

\begingroup
\footnotesize
\setlength{\tabcolsep}{3.5pt}
\renewcommand{\arraystretch}{1.03}
\begin{longtable}{@{}>{\cyrillicfont}p{31mm} p{23mm} r @{\hspace{4mm}} >{\cyrillicfont}p{31mm} p{23mm} r@{}}
\toprule
\normalfont Name & Field & Birth & \normalfont Name & Field & Birth\\
\midrule
\endfirsthead
\toprule
\normalfont Name & Field & Birth & \normalfont Name & Field & Birth\\
\midrule
\endhead
Стеклов & math phys & 1864 & Персидский & diff equ & 1903\\
Крылов & mech & 1879 & Норден & geom & 1904\\
Бернштейн & approx & 1880 & Чудаков & no. theor & 1904\\
Смирнов & anal & 1887 & Моисеев & stab & 1905\\
Делоне & no. theor & 1890 & Райков & groups & 1905\\
Виноградов & no. theor & 1891 & Левин & real \& compl var & 1906\\
Мусхелишвили & elast & 1891 & Тихонов & topol & 1906\\
Привалов & compl var & 1891 & Юшкевич & hist math & 1906\\
Федоров & compl var & 1893 & Векуа, И. Н. & funct anal & 1907\\
Чеботарев & alg & 1894 & Крейн & funct anal & 1907\\
Александров, П. С. & topol & 1896 & Рашевский & diff geom & 1907\\
Гончаров & real var & 1896 & Фаддеев & alg & 1907\\
Лаппо-Данилевский & diff equ & 1896 & Вагнер & diff geom & 1908\\
Гагаев & series & 1897 & Гантмахер & matrices & 1908\\
Гринберг & diff equ & 1900 & Маркушевич & compl var & 1908\\
Немыцкий & diff equ & 1900 & Соболев & funct anal & 1908\\
Андронов & mech & 1901 & Ефимов & geom & 1910\\
Ахиезер & series & 1901 & Келдыш & compl var & 1911\\
Новиков & alg & 1901 & Гнеденко & prob & 1912\\
Петровский & diff equ & 1901 & Черников & alg & 1912\\
Гаврилов & cont sys & 1903 & Левитан & series & 1914\\
Колмогоров & prob & 1903 & Бицадзе & diff equ & 1916\\
Купрадзе & int equ & 1903 & Шилов & funct anal & 1917\\
Марков & logic & 1903 & Погорелов & geom & 1919\\
Виленкин & groups & 1920 & Березанский & funct anal & 1925\\
Стечкин & series & 1920 & Ситников & topol & 1926\\
Граев & groups & 1922 & Постников, М. М. & topol & 1927\\
Шафаревич & alg geom & 1923 & & &\\
\bottomrule
\end{longtable}
\endgroup
